% Auto-generated complete chapter from 10_github_linkedin_cv_strategy.md
\chapter{10 — GitHub, LinkedIn and CV Strategy}
\label{complete:root_012}

\begin{sourcemeta}
\textbf{Fuente:} \href{file:///E:/Laboral/10_github_linkedin_cv_strategy.md}{10\_github\_linkedin\_cv\_strategy.md}\\
\textbf{Seccion:} root\\
\textbf{Estado:} canonical\\
\textbf{Rol:} Modulo principal\\
\textbf{Lineas originales:} 1210\\
\textbf{Ultima modificacion:} 2026-06-11 13:38\\
\textbf{Deduplicacion conservadora:} 0 bloques repetidos omitidos; 0 caracteres repetidos no reimpresos.
\end{sourcemeta}

\section*{Resumen de orientacion}
This file defines how Alexander Woodcock Salomón should structure GitHub, LinkedIn and CV materials for an international job-search strategy focused on:

\section*{Contenido completo depurado}
\addcontentsline{toc}{section}{Contenido completo depurado}




\subsection{0. Purpose of this section}




This file defines how Alexander Woodcock Salomón should structure GitHub, LinkedIn and CV materials for an international job-search strategy focused on:




\begin{enumerate}
\item Real-Time 3D / Interactive Technical Visualization Developer.
\item Unity WebGL / Interactive 3D Developer.
\item Technical Visualization / Digital Twin / Simulation Developer.
\item Unity Technical Artist focused on optimization, runtime interaction and tools.
\item Tools / Pipeline / Python Automation Developer as a controlled secondary route.
\end{enumerate}




The goal is not cosmetic optimization. The goal is to make the public professional profile coherent, defensible and aligned with the strongest evidence: \textbf{TwinSight X500}.




This file uses the source-of-truth decisions from:




\begin{itemize}
\item \texttt{01\_source\_of\_truth\_profile.md}
\item \texttt{02\_market\_role\_fit\_and\_positioning.md}
\item \texttt{07\_portfolio\_strategy\_and\_project\_architecture.md}
\item \texttt{08\_twinsight\_x500\_case\_study.md}
\item \texttt{09\_ara\_ai\_tools\_and\_secondary\_projects.md}
\end{itemize}




It does not create final application-specific CVs for individual job postings. It defines the base strategy, reusable text, variants and quality gates.




\medskip\hrule\medskip




\subsection{1. Core positioning rule}




All public materials must converge on one primary identity:




\begin{quote}
\textbf{Real-Time 3D Developer / Unity Technical Artist focused on interactive technical visualization, CAD-to-realtime optimization and WebGL deployment.}
\end{quote}




Secondary identity:




\begin{quote}
\textbf{Python / AI-assisted tooling developer for research and production workflows.}
\end{quote}




This secondary identity must never overpower the primary identity. A recruiter should understand in less than 10 seconds that the main employable axis is Unity / real-time 3D / technical visualization, not generic full-stack, not pure 3D art, and not AI research.




\medskip\hrule\medskip




\subsection{2. Non-negotiable constraints}




\subsubsection{2.1 Claims that are safe}




The following claims are consistent with the current evidence:




\begin{itemize}
\item Unity / C\# implementation.
\item Unity WebGL interactive 3D prototype.
\item CAD-to-realtime asset optimization.
\item Blender optimization, retopology, UV and baking workflows.
\item Technical UI and runtime interaction systems.
\item Academic usability validation using SUS, NASA-TLX Raw and Think-Aloud.
\item Python automation prototypes.
\item AI-assisted development with human ownership of integration, debugging and technical decisions.
\item English professional working proficiency, if interview performance supports it.
\item Colombia-based remote contractor availability.
\end{itemize}




\subsubsection{2.2 Claims that must be avoided}




Do not use these claims unless future evidence changes:




\begin{itemize}
\item “Senior Technical Artist.”
\item “Senior Unity Developer.”
\item “AI Engineer.”
\item “Machine Learning Engineer.”
\item “Digital Twin Engineer” as a formal job-title identity.
\item “EU work authorized” before Portuguese passport or German residence authorization exists.
\item “Certified” for non-certified courses.
\item “Production-ready AI platform” for ARA.
\item “Commercial shipped Unity product” for TwinSight.
\item “Original drone CAD design.”
\item “Full-stack developer” as the primary headline.
\end{itemize}




\subsubsection{2.3 Correct work authorization language}




Current valid language:




\begin{mdcode}
Bloque de codigo: text
Based in Colombia. Available for remote contractor/B2B work with international teams.
\end{mdcode}




Do not state:




\begin{mdcode}
Bloque de codigo: text
Authorized to work in the EU.
\end{mdcode}




Future language after Portuguese passport is issued:




\begin{mdcode}
Bloque de codigo: text
Portuguese / EU citizen. Authorized to work across the EU/EEA.
\end{mdcode}




Future language if legally resident in Germany through marriage/family reunification and work authorization is issued:




\begin{mdcode}
Bloque de codigo: text
Based in Germany with work authorization.
\end{mdcode}




Until that status exists, it must remain a future mobility possibility, not a current claim.




\medskip\hrule\medskip




\subsection{3. Public profile architecture}




\subsubsection{3.1 Recommended public hierarchy}




The public profile should show this order:




\begin{enumerate}
\item TwinSight X500.
\item ARA Framework, only if cleaned and presented as tooling.
\item Hyperrealistic Blender portrait, only as a technical art / asset pipeline breakdown.
\item Optional ai-news-aggregator, only if functional and documented.
\item Older repositories only if cleaned, contextualized or hidden.
\end{enumerate}




\subsubsection{3.2 What each platform should prove}




\begin{center}
\scriptsize
\begin{longtable}{>{\raggedright\arraybackslash}p{0.455\linewidth} >{\raggedright\arraybackslash}p{0.455\linewidth}}
\toprule
Platform & Main proof \\
\midrule
\endfirsthead
\toprule
Platform & Main proof \\
\midrule
\endhead
CV & Fit, keywords, quantified evidence. \\
LinkedIn & Recruiter discoverability and credibility. \\
GitHub & Technical implementation and documentation. \\
Portfolio & Project quality and visual/technical storytelling. \\
ArtStation & Visual pipeline credibility. \\
\bottomrule
\end{longtable}
\end{center}




The CV should not carry everything. GitHub should not be treated as a storage dump. LinkedIn should not be a biography. Each platform has a distinct job.




\medskip\hrule\medskip




\section{Part A — GitHub Strategy}




\subsection{4. GitHub objective}




GitHub must show that Alexander can structure, document and maintain technical projects, not merely upload code.




For the target roles, GitHub should prove:




\begin{itemize}
\item Unity project organization.
\item C\# system structure.
\item technical documentation quality;
\item reproducible project setup;
\item understanding of asset pipelines;
\item clear ownership and limitations;
\item controlled use of AI assistance;
\item ability to communicate technical tradeoffs.
\end{itemize}




A messy GitHub weakens the profile. A smaller but curated GitHub is better than many incomplete repositories.




\medskip\hrule\medskip




\subsection{5. GitHub profile cleanup}




\subsubsection{5.1 Profile bio}




Recommended GitHub bio:




\begin{mdcode}
Bloque de codigo: text
Real-Time 3D Developer / Unity Technical Artist focused on interactive technical visualization, CAD-to-realtime optimization and WebGL deployment.
\end{mdcode}




Alternative shorter version:




\begin{mdcode}
Bloque de codigo: text
Unity • WebGL • Real-Time 3D • Technical Visualization • CAD-to-Realtime
\end{mdcode}




\subsubsection{5.2 Profile links}




The GitHub profile should link to:




\begin{enumerate}
\item Portfolio website.
\item LinkedIn.
\item ArtStation, only after meaningful content exists.
\item Email or contact method, only if professional and stable.
\end{enumerate}




Do not link to unfinished portfolio branches as the main link unless the page is presentable.




\subsubsection{5.3 Pinned repositories}




Recommended pinned repository order:




\begin{enumerate}
\item \texttt{WebGL-Thesis-Proposal} / \texttt{TwinSight-X500}
\item \texttt{ARA-Framework}, only after cleanup.
\item \texttt{Blender-Technical-Art-Breakdown} or equivalent, only if converted into a portfolio repo.
\item \texttt{ai-news-aggregator}, only if functional.
\item \texttt{Unity-WebGL-Demo} or future smaller polished project.
\item Optional: portfolio website repository.
\end{enumerate}




If GitHub allows only six pinned repositories, do not pin low-signal repos such as tutorial repos, broken experiments or unclean client work.




\medskip\hrule\medskip




\subsection{6. Repository decision matrix}




\begin{center}
\scriptsize
\begin{longtable}{>{\raggedright\arraybackslash}p{0.455\linewidth} >{\raggedright\arraybackslash}p{0.455\linewidth}}
\toprule
Repository & Decision \\
\midrule
\endfirsthead
\toprule
Repository & Decision \\
\midrule
\endhead
\texttt{WebGL-Thesis-Proposal} & Public flagship. Mandatory cleanup. \\
ARA Framework & Conditional public release. \\
ai-news-aggregator & Keep private until functional. \\
FitApp-Free & Keep private or low priority. \\
Hyperrealistic portrait & Publish as portfolio breakdown, not code repo. \\
\texttt{skills-introduction-to-github} & Do not pin. Archive if needed. \\
old client/web repos & Hide unless clean and relevant. \\
\texttt{deus-ex-machina} & Audit before showing. \\
\texttt{delarge95.github.io} & Use only if portfolio-ready. \\
\bottomrule
\end{longtable}
\end{center}




\medskip\hrule\medskip




\subsection{7. TwinSight GitHub README structure}




The TwinSight README must be recruiter-readable first and technically detailed second.




Recommended structure:




\begin{mdcode}
Bloque de codigo: text
# TwinSight X500 — Unity WebGL Technical Visualization Prototype

## Overview
## Live Demo / Video Demo
## Why This Project Exists
## Core Features
## Technical Stack
## CAD-to-Realtime Pipeline
## Runtime Systems
## Optimization Results
## Usability Evaluation
## My Role and Contributions
## AI-Assisted Development Disclosure
## Limitations
## Repository Structure
## Setup Notes
## Screenshots
## Related Links
\end{mdcode}




\subsubsection{7.1 README opening paragraph}




Recommended opening:




\begin{mdcode}
Bloque de codigo: text
TwinSight X500 is a Unity WebGL interactive technical visualization prototype for inspecting the assembly structure of a Holybro X500 V2 drone. It converts CAD/manufacturer geometry and technical documentation into an optimized real-time 3D experience with component selection, exploded view, clipping, technical information panels and multiple visualization modes.
\end{mdcode}




\subsubsection{7.2 Required metrics block}




Include a compact metric block:




\begin{mdcode}
Bloque de codigo: text
Key metrics:
- CAD routes reduced from >6.5M triangles to 95,617 optimized triangles.
- 12 validation participants.
- SUS average: 91.88.
- NASA-TLX Raw: 8.69 for 3D viewer vs 19.89 for 2D support.
- 96 task-condition records.
\end{mdcode}




Add a note that these are academic evaluation metrics, not commercial product KPIs.




\subsubsection{7.3 AI disclosure block}




Recommended wording:




\begin{mdcode}
Bloque de codigo: text
Development note: AI tools were used as coding and documentation assistants during implementation. Final architecture, integration, debugging, Unity scene assembly, technical decisions, asset optimization, evaluation design and project ownership remained my responsibility.
\end{mdcode}




This is direct enough to be honest but avoids weakening the project unnecessarily.




\medskip\hrule\medskip




\subsection{8. ARA GitHub README structure}




ARA must be positioned as tooling, not as a production AI company product.




Recommended title:




\begin{mdcode}
Bloque de codigo: text
ARA Framework — AI-Assisted Research Automation Prototype
\end{mdcode}




Recommended README structure:




\begin{mdcode}
Bloque de codigo: text
# ARA Framework

## Overview
## Problem
## Architecture
## Agent Workflow
## Tech Stack
## Example Output
## Current Status
## Limitations
## Setup
## Roadmap
## AI and API Disclosure
\end{mdcode}




\subsubsection{8.1 ARA positioning statement}




\begin{mdcode}
Bloque de codigo: text
ARA Framework is a Python-based research automation prototype that explores structured academic and technical report generation using LangGraph, LangChain, LLM APIs, retrieval tools, checkpointing and Markdown output workflows.
\end{mdcode}




\subsubsection{8.2 ARA claims to avoid}




Do not describe ARA as:




\begin{itemize}
\item autonomous research scientist;
\item production-grade AI platform;
\item enterprise-ready agent system;
\item replacement for academic review;
\item evidence of ML engineering.
\end{itemize}




\medskip\hrule\medskip




\subsection{9. GitHub quality gates before applying}




Before sending applications, GitHub must satisfy these conditions:




\begin{itemize}
\item Profile bio aligned with target role.
\item TwinSight pinned first.
\item TwinSight README complete.
\item Broken or irrelevant repos unpinned.
\item Private repos hidden if unfinished.
\item At least one demo video linked.
\item At least one screenshot set linked.
\item No secret keys, API keys or credentials.
\item No unprofessional commit messages in visible history.
\item No overclaiming in README files.
\item AI assistance disclosed where relevant.
\end{itemize}




If these conditions are not met, LinkedIn and CV may generate interest but technical reviewers may reject the profile after checking GitHub.




\medskip\hrule\medskip




\section{Part B — LinkedIn Strategy}




\subsection{10. LinkedIn objective}




LinkedIn must make Alexander discoverable for international recruiters and technical leads searching for:




\begin{itemize}
\item Unity Developer;
\item Technical Artist;
\item Real-Time 3D Developer;
\item Interactive 3D Developer;
\item Unity WebGL Developer;
\item Technical Visualization Developer;
\item Digital Twin Visualization;
\item Simulation Developer;
\item CAD-to-realtime;
\item Technical Art;
\item Python automation;
\item AI-assisted tooling.
\end{itemize}




The profile must be written in English as the primary international job-search language.




\medskip\hrule\medskip




\subsection{11. LinkedIn headline options}




\subsubsection{Recommended final headline}




\begin{mdcode}
Bloque de codigo: text
Real-Time 3D Developer / Unity Technical Artist | Unity WebGL, C#, CAD-to-Realtime Optimization, Technical Visualization
\end{mdcode}




This is the best balance between clarity, ATS keywords and defensibility.




\subsubsection{Alternative headlines}




\begin{enumerate}
\item \texttt{Unity Technical Artist / Real-Time 3D Developer | WebGL, C\#, Blender, CAD-to-Realtime Workflows}
\item \texttt{Interactive 3D Developer | Unity WebGL, Technical Visualization, CAD Optimization, C\#}
\item \texttt{Real-Time 3D Developer | Unity, WebGL, Technical UI, CAD-to-Realtime Visualization}
\item \texttt{Unity WebGL Developer | Interactive Technical Visualization, C\#, Blender, Real-Time 3D}
\item \texttt{Technical Artist \& Real-Time 3D Developer | Unity URP, WebGL, Optimization, Visualization Systems}
\item \texttt{Unity Developer for Interactive 3D Visualization | C\#, WebGL, CAD-to-Realtime Pipelines}
\item \texttt{Real-Time Technical Visualization Developer | Unity, WebGL, Blender, C\#}
\item \texttt{Unity Technical Artist | Asset Optimization, Technical UI, WebGL, C\#}
\item \texttt{Interactive 3D / Unity Developer | Technical Visualization, WebGL Deployment, CAD Optimization}
\item \texttt{Technical Visualization Developer | Unity WebGL, Real-Time 3D, C\#, Blender}
\item \texttt{Unity / WebGL Developer | Real-Time 3D, Technical Art, Interactive Visualization}
\item \texttt{CAD-to-Realtime Technical Artist | Unity, Blender, WebGL, Interactive 3D Systems}
\item \texttt{Multimedia Engineer | Unity WebGL, Real-Time 3D, Technical Visualization, Python Automation}
\item \texttt{Unity Technical Artist \& Tools Developer | Real-Time 3D, WebGL, Python Automation}
\end{enumerate}

15. \texttt{Technical Artist in Training | Unity WebGL, C\#, Real-Time 3D, Technical Visualization}




Avoid “Senior” in the headline. Avoid leading with Python or AI unless targeting a Python-specific role.




\medskip\hrule\medskip




\subsection{12. LinkedIn About section}




Recommended About:




\begin{mdcode}
Bloque de codigo: text
I am a Multimedia Engineering student completing my thesis defense, focused on real-time 3D, Unity WebGL and interactive technical visualization.

My flagship project, TwinSight X500, is a Unity WebGL prototype for inspecting the assembly structure of a Holybro X500 V2 drone. The project transforms CAD/manufacturer geometry and technical documentation into an optimized real-time 3D experience with component selection, exploded view, clipping, technical information panels and multiple visualization modes.

My strongest areas are Unity, C#, WebGL deployment, CAD-to-realtime optimization, Blender asset preparation, technical UI, real-time interaction systems and technical documentation. I also build Python-based automation prototypes, including AI-assisted research tooling with LangGraph and LLM workflows.

I am currently based in Colombia and available for remote contractor/B2B opportunities with international teams. I am primarily targeting roles in real-time 3D, Unity technical art, interactive 3D visualization, technical visualization, simulation, XR and digital twin visualization.

Core stack: Unity, C#, WebGL, URP, UI Toolkit, Shader Graph, Blender, Python, Git, technical documentation.
\end{mdcode}




\subsubsection{12.1 About section rationale}




This version does four things:




\begin{enumerate}
\item Leads with current role identity.
\item Makes TwinSight visible immediately.
\item Includes relevant keywords without keyword stuffing.
\item States Colombia remote availability without pretending to have EU work authorization.
\end{enumerate}




\medskip\hrule\medskip




\subsection{13. LinkedIn Experience section}




\subsubsection{13.1 Recommended current role title}




Use:




\begin{mdcode}
Bloque de codigo: text
Independent Technical Artist & Real-Time 3D Developer
\end{mdcode}




Alternative:




\begin{mdcode}
Bloque de codigo: text
Independent Unity / Real-Time 3D Developer
\end{mdcode}




Avoid:




\begin{mdcode}
Bloque de codigo: text
Senior Technical Artist
AI Engineer
Full-Stack Developer
Founder
\end{mdcode}




\subsubsection{13.2 Recommended experience entry}




\begin{mdcode}
Bloque de codigo: text
Independent Technical Artist & Real-Time 3D Developer
Freelance / Portfolio Projects
Mar 2024 – Present
Colombia · Remote / Hybrid

- Developed TwinSight X500, a Unity WebGL interactive technical visualization prototype for inspecting the assembly structure of a Holybro X500 V2 drone.
- Built runtime interaction systems including component selection, exploded view, clipping/cross-section controls, visualization modes and technical information panels.
- Optimized CAD/manufacturer geometry for real-time WebGL deployment, reducing high-density CAD routes from more than 6.5M triangles to 95,617 optimized triangles.
- Prepared assets through Blender-based cleanup, retopology, UV and high-to-low workflow planning.
- Designed and documented an academic validation process using SUS, NASA-TLX Raw and Think-Aloud methods.
- Explored Python automation and AI-assisted research tooling through ARA Framework, a LangGraph-based academic research automation prototype.
\end{mdcode}




\subsubsection{13.3 Notes}




The experience entry should not read like traditional employment if it was primarily academic/portfolio work. “Freelance / Portfolio Projects” is clearer than pretending it was formal industry employment.




\medskip\hrule\medskip




\subsection{14. LinkedIn Education section}




\subsubsection{14.1 UNAD entry}




Recommended:




\begin{mdcode}
Bloque de codigo: text
Universidad Nacional Abierta y a Distancia — UNAD
Bachelor of Engineering, Multimedia Engineering
Aug 2021 – Jul 2026

Relevant areas: computer graphics, software engineering, mathematics, physics, real-time 3D, Unity development, technical visualization and interactive systems.

Thesis: TwinSight X500 — Unity WebGL interactive technical visualization prototype for drone assembly inspection.
\end{mdcode}




Important correction: do not describe ARA Framework as the thesis unless formally true.




\subsubsection{14.2 UNAL entry}




Recommended:




\begin{mdcode}
Bloque de codigo: text
Universidad Nacional de Colombia — UNAL
Advanced Coursework in Electronic Engineering
Mar 2013 – Nov 2018

Completed advanced coursework in electronic engineering, including mathematics, physics, circuits, systems, simulation and algorithmic problem-solving. This background supports technical visualization, engineering-oriented 3D work and spatial/analytical reasoning.
\end{mdcode}




Alternative shorter wording:




\begin{mdcode}
Bloque de codigo: text
Coursework toward B.Eng. in Electronic Engineering
\end{mdcode}




Avoid saying “graduated” from UNAL if no degree was obtained.




\medskip\hrule\medskip




\subsection{15. LinkedIn Projects section}




\subsubsection{15.1 TwinSight X500}




\begin{mdcode}
Bloque de codigo: text
TwinSight X500 — Unity WebGL Technical Visualization Prototype

Interactive 3D technical visualization system for inspecting the assembly structure of a Holybro X500 V2 drone. Built with Unity WebGL, C#, URP, UI Toolkit and Blender-based CAD-to-realtime optimization.

Features include component selection, exploded view, clipping/cross-section interaction, technical information panels and multiple visualization modes. The project reduced high-density CAD routes from more than 6.5M triangles to 95,617 optimized triangles and was evaluated using SUS, NASA-TLX Raw and Think-Aloud methods.
\end{mdcode}




\subsubsection{15.2 ARA Framework}




\begin{mdcode}
Bloque de codigo: text
ARA Framework — AI-Assisted Research Automation Prototype

Python-based prototype for structured academic and technical research automation using LangGraph, LangChain, LLM APIs, Redis, Semantic Scholar integration, Playwright-based retrieval and Markdown report generation.

Positioned as a tooling and automation project, not a production AI platform.
\end{mdcode}




\subsubsection{15.3 Blender portrait}




\begin{mdcode}
Bloque de codigo: text
Hyperrealistic Blender Portrait — Technical Art Breakdown

Blender character/portrait study focused on topology, materials, grooming, lighting and rendering. Used as evidence of 3D production literacy and asset pipeline understanding.
\end{mdcode}




Only add this after there is a public ArtStation or portfolio breakdown.




\medskip\hrule\medskip




\subsection{16. LinkedIn Featured section}




Recommended order:




\begin{enumerate}
\item TwinSight case study page.
\item TwinSight demo video.
\item TwinSight GitHub repository.
\item Portfolio homepage.
\item ARA Framework, only after cleanup.
\item ArtStation breakdown, only after cleanup.
\end{enumerate}




Do not feature unfinished branches, old plans or weak repositories.




\medskip\hrule\medskip




\subsection{17. LinkedIn skills order}




Recommended top skills:




\begin{enumerate}
\item Unity
\item C\#
\item Real-Time 3D
\item Technical Art
\item WebGL
\item 3D Optimization
\item Blender
\item Computer Graphics
\item Technical Visualization
\item Shader Graph
\item UI Toolkit
\item Python
\item Git
\item Technical Documentation
\item CAD-to-Realtime Workflows
\item Asset Optimization
\item Interactive 3D
\item XR Foundations
\item Research
\item Software Engineering
\end{enumerate}




If LinkedIn does not allow custom skills such as “CAD-to-Realtime Workflows,” use the closest available equivalent and include the phrase in About, Experience and Projects.




\medskip\hrule\medskip




\subsection{18. LinkedIn activity strategy}




Post only when there is technical substance. Do not post motivational content or generic job-search posts.




Recommended post types:




\begin{enumerate}
\item TwinSight technical breakdown.
\item CAD optimization before/after.
\item Unity WebGL performance note.
\item Clipping/cross-section feature demo.
\item Technical UI breakdown.
\item Lessons from converting CAD to real-time assets.
\item Short demo video with role-targeted caption.
\item ARA architecture note, only after cleanup.
\item Blender topology/material breakdown.
\end{enumerate}




Suggested cadence:




\begin{itemize}
\item 1 strong post per week for 8 weeks.
\item 3–5 thoughtful comments per week on posts from Unity, XR, technical art, simulation and visualization professionals.
\item 5–10 recruiter or technical lead connection requests per week after profile cleanup.
\end{itemize}




Avoid daily posting unless the quality is high.




\medskip\hrule\medskip




\subsection{19. LinkedIn connection scripts}




\subsubsection{19.1 Recruiter connection}




\begin{mdcode}
Bloque de codigo: text
Hi [Name], I’m a Colombia-based Real-Time 3D / Unity developer focused on interactive technical visualization, Unity WebGL and CAD-to-realtime optimization. I’m currently preparing my thesis project, TwinSight X500, as a portfolio case study and I’m exploring remote international roles. Open to connecting.
\end{mdcode}




\subsubsection{19.2 Technical lead connection}




\begin{mdcode}
Bloque de codigo: text
Hi [Name], I’m working on Unity WebGL technical visualization and CAD-to-realtime optimization. My flagship project is TwinSight X500, an interactive 3D drone assembly inspection prototype. I’d like to follow your work in real-time 3D / visualization.
\end{mdcode}




\subsubsection{19.3 Referral request}




\begin{mdcode}
Bloque de codigo: text
Hi [Name], I saw the [Role] opening at [Company]. My strongest fit is Unity/WebGL technical visualization, CAD-to-realtime optimization and interactive 3D systems through my TwinSight X500 project. Would you be open to pointing me to the right recruiter or sharing whether this role accepts remote contractors from LATAM?
\end{mdcode}




\medskip\hrule\medskip




\section{Part C — CV Strategy}




\subsection{20. CV objective}




The CV must function as an ATS-readable proof document. It should not be a life story.




The CV must answer:




\begin{enumerate}
\item What role is he applying for?
\item What is the strongest evidence?
\item What technologies can he defend?
\item Is he available remotely from Colombia?
\item What are the quantified results?
\item What should the recruiter click next?
\end{enumerate}




The default CV should be one page. A two-page version is acceptable only for highly technical roles where project detail is requested.




\medskip\hrule\medskip




\subsection{21. Base CV structure}




Recommended order:




\begin{mdcode}
Bloque de codigo: text
Name + Title
Location / Remote availability / Links
Professional Summary
Technical Skills
Selected Projects
Experience
Education
Additional Training
Languages
\end{mdcode}




If applying to roles that heavily value formal experience, place Experience before Projects. For Alexander’s current stage, \textbf{Selected Projects should come before Experience} because TwinSight is the strongest proof.




\medskip\hrule\medskip




\subsection{22. Base CV header}




\begin{mdcode}
Bloque de codigo: text
Alexander Woodcock Salomón
Real-Time 3D Developer / Unity Technical Artist
Colombia · Remote Contractor/B2B · English C1 · Spanish Native
LinkedIn: [link] · GitHub: [link] · Portfolio: [link] · ArtStation: [link if ready]
\end{mdcode}




Do not include full address. City/country is enough.




\medskip\hrule\medskip




\subsection{23. Base CV professional summary}




\begin{mdcode}
Bloque de codigo: text
Real-Time 3D Developer / Unity Technical Artist focused on Unity WebGL, C#, CAD-to-realtime optimization and interactive technical visualization. Built TwinSight X500, a Unity WebGL drone assembly visualization prototype with component selection, exploded view, clipping, technical UI and multiple visualization modes. Strong background in Blender asset optimization, technical documentation, computer graphics fundamentals and Python automation prototypes. Based in Colombia and available for remote contractor/B2B work with international teams.
\end{mdcode}




This is the default version. It can be adapted by role family.




\medskip\hrule\medskip




\subsection{24. Technical skills section}




Recommended structure:




\begin{mdcode}
Bloque de codigo: text
Real-Time 3D: Unity, C#, Unity WebGL, URP, UI Toolkit, Shader Graph, runtime interaction systems
3D Pipeline: Blender, retopology, UVs, baking, CAD-to-realtime optimization, asset cleanup
Visualization: technical visualization, exploded views, clipping/cross-section interaction, technical UI
Programming / Tools: Python, Git/GitHub, LangGraph, LangChain, FastAPI basics, React basics
Documentation / Evaluation: technical documentation, SUS, NASA-TLX Raw, Think-Aloud, academic validation
\end{mdcode}




Only include secondary technologies like React, PostgreSQL, Redis or FastAPI if the target role values them and there is a project that demonstrates them.




\medskip\hrule\medskip




\subsection{25. Selected projects — CV entries}




\subsubsection{25.1 TwinSight X500}




\begin{mdcode}
Bloque de codigo: text
TwinSight X500 — Unity WebGL Technical Visualization Prototype
Academic thesis project · Unity, C#, WebGL, URP, UI Toolkit, Blender

- Built an interactive Unity WebGL prototype for inspecting the assembly structure of a Holybro X500 V2 drone through component selection, exploded view, clipping and technical information panels.
- Converted CAD/manufacturer geometry into optimized real-time assets, reducing high-density CAD routes from more than 6.5M triangles to 95,617 optimized triangles.
- Implemented multiple visualization modes including realistic, X-Ray, blueprint, solid color, wireframe, ghosted and thermal-style views.
- Designed technical UI and runtime interaction systems for component inspection and spatial relationship understanding.
- Evaluated the prototype with 12 participants using SUS, NASA-TLX Raw and Think-Aloud methods; recorded 96 task-condition observations.
\end{mdcode}




\subsubsection{25.2 ARA Framework}




\begin{mdcode}
Bloque de codigo: text
ARA Framework — AI-Assisted Research Automation Prototype
Python, LangGraph, LangChain, Redis, Semantic Scholar API, Playwright MCP

- Developed a Python prototype for structured academic and technical report generation using a multi-agent workflow with LangGraph and LLM APIs.
- Designed agent roles for niche analysis, literature research, technical architecture, implementation planning and content synthesis.
- Integrated retrieval, checkpointing, budget management and Markdown report generation workflows.
- Positioned as a research automation prototype and tooling experiment, not as a production AI platform.
\end{mdcode}




Use ARA only in CV variants where Python/tooling helps the target job.




\subsubsection{25.3 Blender portrait}




\begin{mdcode}
Bloque de codigo: text
Hyperrealistic Blender Portrait — Technical Art Study
Blender, topology, materials, grooming, lighting, rendering

- Completed a Blender character/portrait study focused on topology, material setup, grooming, lighting and rendering.
- Structured the project as evidence of 3D production literacy and technical art pipeline understanding.
\end{mdcode}




Only include if the portfolio/ArtStation page exists.




\medskip\hrule\medskip




\subsection{26. Experience section}




Recommended entry:




\begin{mdcode}
Bloque de codigo: text
Independent Technical Artist & Real-Time 3D Developer
Freelance / Portfolio Projects · Colombia · Mar 2024 – Present

- Developed Unity, WebGL, Python and 3D portfolio projects focused on interactive visualization, technical documentation and production workflows.
- Built and documented TwinSight X500 as a thesis-based technical visualization prototype.
- Explored AI-assisted research automation through Python and LangGraph-based tooling.
- Completed freelance web and digital projects for small clients, strengthening client communication, delivery discipline and technical troubleshooting.
\end{mdcode}




If formal freelance projects are weak or unrelated, keep this section short and let Projects dominate.




\medskip\hrule\medskip




\subsection{27. Education section}




\subsubsection{27.1 UNAD}




\begin{mdcode}
Bloque de codigo: text
Universidad Nacional Abierta y a Distancia — UNAD
B.Eng. Multimedia Engineering · Expected 2026
Thesis: TwinSight X500 — Unity WebGL interactive technical visualization prototype for drone assembly inspection.
Relevant areas: computer graphics, software engineering, mathematics, physics, real-time 3D, Unity development and technical visualization.
\end{mdcode}




\subsubsection{27.2 UNAL}




\begin{mdcode}
Bloque de codigo: text
Universidad Nacional de Colombia — UNAL
Advanced Coursework in Electronic Engineering · 2013–2018
Completed advanced coursework in mathematics, physics, electronics, circuits, systems, simulation and algorithmic problem-solving.
\end{mdcode}




This wording is honest and internationally understandable.




\medskip\hrule\medskip




\subsection{28. Additional training section}




Because many courses are non-certified, they should not be presented as formal credentials.




Recommended heading:




\begin{mdcode}
Bloque de codigo: text
Additional Training and Self-Directed Study
\end{mdcode}




Recommended format:




\begin{mdcode}
Bloque de codigo: text
- CG Cookie HUMAN — High-Fidelity Character Pipeline & Topology; completed as self-directed technical art training.
- Alive! — 3D Animation Mechanics & Dynamic Motion; completed as self-directed animation study.
- Rebelway / Houdini / Unreal FX coursework; used for portfolio-directed skill development. [Only after meaningful progress.]
\end{mdcode}




Do not write:




\begin{mdcode}
Bloque de codigo: text
Certified in CG Cookie HUMAN
Certified Rebelway FX Artist
\end{mdcode}




unless a formal certificate exists and can be verified.




\medskip\hrule\medskip




\subsection{29. Languages section}




\begin{mdcode}
Bloque de codigo: text
Spanish — Native
English — Full professional proficiency / C1 self-assessed
German — Beginner / A1 target if Germany route becomes active
Portuguese — Future strategic language for Portugal/EU mobility; not current work language unless studied
\end{mdcode}




For a one-page CV, include only Spanish and English unless German or Portuguese becomes relevant to the job.




\medskip\hrule\medskip




\section{Part D — CV Variants}




\subsection{30. Variant 1 — Unity Technical Artist / Real-Time 3D Developer}




\subsubsection{Summary}




\begin{mdcode}
Bloque de codigo: text
Unity Technical Artist / Real-Time 3D Developer focused on interactive technical visualization, Unity WebGL, asset optimization and technical UI. Built TwinSight X500, a drone assembly visualization prototype combining CAD-to-realtime optimization, runtime interaction systems, clipping, exploded view and usability validation.
\end{mdcode}




\subsubsection{Emphasize}




\begin{itemize}
\item Unity.
\item C\#.
\item WebGL.
\item Blender optimization.
\item Shader Graph.
\item Technical UI.
\item CAD-to-realtime.
\end{itemize}




\subsubsection{Downplay}




\begin{itemize}
\item ARA.
\item Full-stack.
\item FitApp-Free.
\end{itemize}




\medskip\hrule\medskip




\subsection{31. Variant 2 — Unity WebGL / Interactive 3D Developer}




\subsubsection{Summary}




\begin{mdcode}
Bloque de codigo: text
Unity WebGL / Interactive 3D Developer focused on browser-based real-time visualization, C# interaction systems and optimized 3D assets. Built TwinSight X500, a Unity WebGL prototype for technical drone assembly inspection with selection, exploded view, clipping and technical information panels.
\end{mdcode}




\subsubsection{Emphasize}




\begin{itemize}
\item WebGL deployment.
\item runtime performance constraints;
\item interactive UI;
\item browser-accessible 3D;
\item mobile-first considerations.
\end{itemize}




\subsubsection{Add if true}




\begin{itemize}
\item live demo URL;
\item build size;
\item load time;
\item FPS / frame-time data.
\end{itemize}




\medskip\hrule\medskip




\subsection{32. Variant 3 — Technical Visualization / Digital Twin / Simulation}




\subsubsection{Summary}




\begin{mdcode}
Bloque de codigo: text
Technical Visualization Developer focused on real-time 3D systems for explaining complex assemblies and spatial relationships. Built TwinSight X500, a Unity WebGL drone assembly inspection prototype combining CAD-to-realtime optimization, technical UI, exploded view, clipping and structured usability evaluation.
\end{mdcode}




\subsubsection{Emphasize}




\begin{itemize}
\item engineering-oriented visualization;
\item CAD/documentation conversion;
\item technical communication;
\item component hierarchy;
\item spatial relationships;
\item usability evaluation.
\end{itemize}




\subsubsection{Avoid}




Do not claim live sensor data, simulation physics, IoT integration or true digital twin functionality unless added later.




\medskip\hrule\medskip




\subsection{33. Variant 4 — Tools / Pipeline / Python Automation}




\subsubsection{Summary}




\begin{mdcode}
Bloque de codigo: text
Tools-oriented Unity / Python developer focused on technical workflows, automation and structured documentation. Built TwinSight X500 as a real-time 3D visualization prototype and ARA Framework as a Python/LangGraph research automation prototype for structured academic and technical report generation.
\end{mdcode}




\subsubsection{Emphasize}




\begin{itemize}
\item Python;
\item LangGraph;
\item workflow automation;
\item documentation;
\item reproducibility;
\item tooling mindset;
\item technical process design.
\end{itemize}




\subsubsection{Downplay}




\begin{itemize}
\item pure art;
\item game-only framing;
\item incomplete full-stack projects.
\end{itemize}




\medskip\hrule\medskip




\section{Part E — Keyword Strategy}




\subsection{34. Primary ATS keyword cluster}




Use these consistently across CV, LinkedIn, GitHub and portfolio:




\begin{mdcode}
Bloque de codigo: text
Unity
C#
Unity WebGL
Real-Time 3D
Technical Artist
Interactive 3D
Technical Visualization
CAD-to-Realtime
3D Optimization
Blender
URP
UI Toolkit
Shader Graph
Runtime Interaction Systems
Technical UI
Exploded View
Clipping / Cross-Section
WebGL Deployment
Computer Graphics
Technical Documentation
\end{mdcode}




\subsection{35. Secondary keyword cluster}




Use only where relevant:




\begin{mdcode}
Bloque de codigo: text
Python
LangGraph
LangChain
LLM Applications
AI-Assisted Tooling
Research Automation
FastAPI
React
PostgreSQL
Redis
Playwright
Semantic Scholar API
Markdown Automation
\end{mdcode}




\subsection{36. Industry keyword cluster}




Use selectively by job type:




\begin{mdcode}
Bloque de codigo: text
Digital Twin Visualization
Simulation
XR
AR/VR
Product Visualization
Industrial Visualization
Technical Training
Interactive Documentation
Drone Assembly
Engineering Visualization
Manufacturing Visualization
\end{mdcode}




\subsection{37. Spanish keyword equivalents}




Use in LATAM searches or Spanish CV variants:




\begin{mdcode}
Bloque de codigo: text
Desarrollador Unity
Artista Técnico
Desarrollador 3D en Tiempo Real
Visualización Técnica
Optimización 3D
WebGL
Gemelo Digital
Simulación
XR / Realidad Extendida
Automatización con Python
Documentación Técnica
\end{mdcode}




\medskip\hrule\medskip




\section{Part F — Quality Control}




\subsection{38. Consistency checklist}




Before applying, verify that:




\begin{itemize}
\item LinkedIn headline matches CV title.
\item CV title matches portfolio title.
\item GitHub pinned repo supports the headline.
\item TwinSight appears first everywhere.
\item ARA appears secondary everywhere.
\item No document says the thesis is ARA if the thesis is TwinSight.
\item No document says thesis is still in early development.
\item No course without certificate is labeled “certified.”
\item No EU work authorization is claimed prematurely.
\item GitHub does not expose weak repos above strong ones.
\item LinkedIn Featured does not show unfinished pages.
\item CV does not look like a generic full-stack profile.
\end{itemize}




\medskip\hrule\medskip




\subsection{39. Overclaiming risk map}




\subsubsection{High-risk claims}




\begin{mdcode}
Bloque de codigo: text
Senior Technical Artist
AI Engineer
Digital Twin Engineer
Production AI Platform
EU Work Authorized
Certified Houdini / Unity / Rebelway without proof
\end{mdcode}




These should be removed.




\subsubsection{Medium-risk claims}




\begin{mdcode}
Bloque de codigo: text
Technical Artist
Tools Developer
Simulation Developer
XR Developer
Full-stack Developer
\end{mdcode}




These are acceptable only when supported by project evidence.




\subsubsection{Low-risk claims}




\begin{mdcode}
Bloque de codigo: text
Unity WebGL
Real-Time 3D
Interactive Technical Visualization
CAD-to-Realtime Optimization
Python Automation Prototype
Technical Documentation
\end{mdcode}




These align with current evidence.




\medskip\hrule\medskip




\subsection{40. Recruiter-facing work availability}




\subsubsection{Current version}




\begin{mdcode}
Bloque de codigo: text
Based in Colombia. Available for remote contractor/B2B roles with international teams.
\end{mdcode}




\subsubsection{For U.S. companies}




\begin{mdcode}
Bloque de codigo: text
I am based in Colombia and available for remote contractor/B2B collaboration. I do not currently require U.S. payroll employment.
\end{mdcode}




\subsubsection{For EU companies before Portuguese passport}




\begin{mdcode}
Bloque de codigo: text
I am currently based in Colombia and available for remote contractor/B2B work. I do not currently hold EU work authorization, but I expect Portuguese citizenship/passport in approximately two years.
\end{mdcode}




Use the Portuguese citizenship line only if strategically relevant and only after the conversation reaches relocation/work authorization.




\subsubsection{After Portuguese passport}




\begin{mdcode}
Bloque de codigo: text
Portuguese / EU citizen. Authorized to work across the EU/EEA.
\end{mdcode}




\subsubsection{If Germany residence/work authorization through marriage exists}




\begin{mdcode}
Bloque de codigo: text
Based in Germany with valid residence and work authorization.
\end{mdcode}




\medskip\hrule\medskip




\subsection{41. Recommended immediate edits}




\subsubsection{GitHub}




\begin{enumerate}
\item Rewrite GitHub bio.
\item Pin TwinSight first.
\item Rewrite TwinSight README.
\item Hide/unpin tutorial and weak repos.
\item Add demo video and screenshots.
\item Add AI-assistance disclosure.
\item Add clear limitations section.
\end{enumerate}




\subsubsection{LinkedIn}




\begin{enumerate}
\item Replace headline.
\item Rewrite About.
\item Correct UNAD thesis description.
\item Reframe UNAL as advanced coursework.
\item Add TwinSight as project.
\item Reorder skills.
\item Add Featured links only after assets are ready.
\item Set Open to Work carefully.
\end{enumerate}




\subsubsection{CV}




\begin{enumerate}
\item Create one-page base CV.
\item Create four variants.
\item Use TwinSight metrics.
\item Remove unsupported claims.
\item Include Colombia remote contractor availability.
\item Add portfolio and GitHub links.
\item Keep courses in Additional Training, not Certifications.
\end{enumerate}




\medskip\hrule\medskip




\subsection{42. Final recommended profile package}




The minimum public package before serious applications:




\begin{mdcode}
Bloque de codigo: text
1. One-page CV: Unity / Real-Time 3D variant.
2. LinkedIn rewritten and corrected.
3. GitHub profile cleaned.
4. TwinSight README complete.
5. TwinSight case-study page live.
6. 60–120 second TwinSight demo video.
7. 8–12 high-quality screenshots.
8. ARA hidden or cleaned as secondary.
9. ArtStation updated only if Blender portrait has a technical breakdown.
\end{mdcode}




Applications before this package is ready may still work, but conversion probability will be lower.




\medskip\hrule\medskip




\subsection{43. Open questions before final CV writing}




These must be answered before producing the final CV files:




\begin{enumerate}
\item Is there a live TwinSight WebGL demo URL?
\item Is the GitHub repository clean enough to show now?
\item Are there screenshots ready?
\item Is the demo video recorded?
\item Which metrics are final and defensible after thesis defense?
\item Which courses from the Excel were actually completed?
\item Does ARA run reliably enough to show publicly?
\item Does ai-news-aggregator currently work?
\item Should the CV mention freelance web work or omit it?
\item What email/contact method should appear on public materials?
\end{enumerate}




\medskip\hrule\medskip




\subsection{44. Decision for next document}




This section defines the public profile strategy. The next section should move from profile assets to external market execution:




\begin{mdcode}
Bloque de codigo: text
11_company_targets_job_boards_recruiters.md
\end{mdcode}




That next file should use external research and structured data collection. It requires company lists, job boards, recruiters, communities, role-title searches and links.




Recommended tool use for the next section:




\begin{itemize}
\item Deep Research: yes.
\item Agent mode: yes, if available.
\end{itemize}



